\section{Introduction}
Multi-turn interaction is rapidly becoming the dominant mode in which language models are deployed, yet recent work has documented substantial 
performance degradation in this regime: \cite{lanban2025lost} report an average accuracy drop of 39\% across six generation task---including mathematics---
over 200{,}000 simulated conversations with 15 top open- and closed-weight LLMs, characterizing this as a ``lost-in-conversation'' failure. We study a specific and
under-examined instance of this broader phenomenon: \emph{multi-turn mathematical reasoning with explicit numerical dependencies}, in which each sub-question is fully
specified, but its solution depends on the numerical result of the preceding turn. This setting captures a wide range of practical quantitative tasks---financial calculations,
multi-step engineering analyses, scientific computations---that decompose naturally into chained sub-problems. It is particularly acute for small language models 
(SLMs, $\sim $1B--10B parameters), which are increasingly deployed under cost and latency constraints but whose limited attention capacity leaves them vulnerable to errors
that compound across turns.

The central technical question in this regime is how mathematical intermediate state should be maintained across turns: which equations, constraints, and partial results from
prior reasoning must remain accessible to subsequent turns, and in what representational form. Current approaches uniformly store this state as unstructured text. Prompt 
compression methods such as LLMLingua~\citep{jiang2023llmlingua} operate at the token level, agnostic to which tokens encode mathematical constraints. Reasoning-summarization 
methods such as Reflexion~\citep{shinn2023reflexion} and InftyThink~\citep{INFTYTHINK_CITE} condense prior reasoning into natural-language narratives in which numerical constraints
stated in passing receive no privileged representation. Raw context accumulation simply defers the problem until attention can no longer locate the relevant constraint---a fragility
consistent with \citet{liu2024lost}, who show that language models do not robustly use information in long input contexts, with utilization varying sharply by position. We argue 
these approaches share an assumption---that reasoning state should be flat text---that is poorly matched to mathematical reasoning, where the state needed across turns is intrinsically
typed: invariants, active variables, binding equations and inequalities, invalidated paths, and the remaining reasoning gap, each with a distinct role in subsequent inference. 

We propose \emph{Sparse Logic State Compression} (SLSC), a twin-agent framework in which a reasoning agent $\mathcal{A}$ generates natural-language chains of thought while a curator
agent $\mathcal{B}$ periodically distills accumulated context into a structured state object consisting of atomic mathematical propositions. The reasoning agent's working context
is then replaced by this object, giving each constraint a persistent and addressable location rather than leaving it to be recovered through attention over an ever-growing context.
Figure~\ref{fig:comparison} contrasts the two representations.

\begin{figure*}[!t]
\centering

\noindent
\begin{minipage}[t]{0.48\textwidth}
\begin{mdframed}[
    frametitle={\textbf{Baseline: Natural-Language Summary}},
    frametitlerule=true,
    frametitlebackgroundcolor=gray!15,
    backgroundcolor=gray!3,
    innertopmargin=6pt, innerbottommargin=6pt
]
\scriptsize\textit{Flat text; intermediate constraints embedded in prose.}\\[4pt]
\small
\textit{Context passed to Agent $\mathcal{A}$ at Turn 3:}\\[4pt]
``To approach this problem, I first considered the divisors of $n+7$ and $2n+1$. By examining the linear combination $2(n+7) - (2n+1)$, I observed that this simplifies to 13, which means \colorbox{yellow!40}{any common divisor $d$ must divide 13}. Since 13 is prime, the possible values are $d=1$ or $d=13$. I then proceeded to analyze the case $d=13$, which requires $n+7 \equiv 0 \pmod{13}$, leading to $n \equiv 6 \pmod{13}$. Combined with the bounds, I narrowed down to candidates $n = 6, 19, 32, \ldots$, verifying each against the secondary constraint $\ldots$''\\[6pt]
\hrule\vspace{4pt}
\scriptsize $\sim$185 tokens $\;\bullet\;$ no addressable structure $\;\bullet\;$ key fact mid-sentence
\end{mdframed}
\end{minipage}
\hfill
\begin{minipage}[t]{0.48\textwidth}
\begin{mdframed}[
    frametitle={\textbf{SLSC: Sparse Logic State}},
    frametitlerule=true,
    frametitlebackgroundcolor=blue!15,
    backgroundcolor=blue!3,
    innertopmargin=6pt, innerbottommargin=6pt
]
\scriptsize\textit{Atomic propositions; each independently addressable.}\\[4pt]
\small
\textit{Context passed to Agent $\mathcal{A}$ at Turn 3:}\\[4pt]
\texttt{1.} $d$ divides $n + 7$\\[2pt]
\texttt{2.} $d$ divides $2n + 1$\\[2pt]
\texttt{3.} $2(n+7) - (2n+1) = 13$\\[2pt]
\texttt{4.} \colorbox{yellow!40}{$d$ divides $13$}\\[2pt]
\texttt{5.} $d \in \{1, 13\}$\\[2pt]
\texttt{6.} If $d = 13$: $n \equiv 6 \pmod{13}$\\[2pt]
\texttt{7.} Candidates: $n \in \{6, 19, 32, \ldots\}$\\[6pt]
\hrule\vspace{4pt}
\scriptsize $\sim$55 tokens $\;\bullet\;$ each fact retrievable $\;\bullet\;$ persistent slot per constraint
\end{mdframed}
\end{minipage}

\caption{Comparison of context representations at Turn 3 of a multi-turn mathematical reasoning task. \textbf{Left}: Natural-language summary embeds constraints in prose, requiring attention over unstructured text to locate key facts (highlighted). \textbf{Right}: SLSC maintains atomic propositions in numbered slots, giving each constraint a persistent and addressable location. SLSC reduces token count by $\sim$70\% while improving constraint retrievability.}
\label{fig:comparison}
\end{figure*}

To evaluate SLSC, we construct a 500-problem multi-turn mathematical reasoning dataset with explicit numerical-dependency annotation by decomposing problems from MATH-500
~\citep{lightman2023math500} into chains of dependent sub-questions. We compare three context-maintenance strategies---No Compression (raw accumulation), Natural-Language Summary(a
Reflexion-style condensed narrative), and SLSC---paired with reasoner models spanning four size from the Qwen3.5 family~\citep{} (0.8B, 2B, 4B, 9B), with a fixed \texttt{gemma-3-27b-it}
~\citep{} curator. Across this grid, SLSC improves final-answer accuracy by [Z\%] over No Compression and [Z'\%] over the natural-language summary baseline while reducing per-turn
token consumption by [W\%]. Critically, the accuracy gap between SLSC and No Compression [widens / narrows] as reasoner size [decrease / increase], indicating that the limited attention
capacity of SLMs make them disproportionately reliant on structured state maintenance for chained mathematical reasoning.

We make four contributions: 
\begin{enumerate}[noitemsep, topsep=2pt]
    \item We construct, to our knowledge, the first multi-turn mathematical reasoning dataset with explicit numerical-dependency structure, complementing existing multi-turn benchmarks that 
    focus on underspecification rather than numerical chaining.
    \item We document the severity of context-accumulation failure in multi-turn mathematical reasoning at SLM scale, and show that the benefit of structured state maintenance scales inversely
    with reasoner size---situating this as a distinctly SLM-scale bottleneck and extending the lost-in-conversation and long-context utilization literature to numerically chained settings.
    \item We propose SLSC, a twin-agent state-compression framework, and show that it improves multi-turn final-answer accuracy by [Z\%] and reduce per-turn token consumption by [W\%] over raw-
    context accumulation, with consistent gains across all four evaluated SLM reasoner sizes.
    \item Contrary to our initial hypothesis, we find through ablation that the choice of state representation syntax---Lean-style formal language versus a simplified structured schema---does not
    significantly affect performance: it is the \emph{granularity} of slotted state, not the \emph{formality} of its syntax, that drives the gains. This has practical implications at SLM scale,
    where simplified schema are substantially more reliable to generate than formal-language ones.
\end{enumerate}